\section{Multiple arithmetic zeros cannot remain tangent through the second Puiseux jet}
\label{sec:analytic-multiple-zero-bending}

This section treats zeros of multiplicity $m\ge2$. The simple
case has an additional deformation term, displayed at the end;
no conclusion about simple-zero transversality is inferred.
All statements at a specified critical-line zero are
unconditional. Assumption A makes them apply to every multiple
nontrivial zero of $\zeta$ or a primitive Dirichlet $L$-function.

\begin{lemma}[A positive logarithmic derivative at one horizontal unit]
\label{lem:analytic-reflected-logderivative}
Let $\Lambda$ be an entire order-one completion with zeros in
$0\le\Re s\le1$, infinitely many zeros, and reflection
\[
 \Lambda(s)=\epsilon\overline{\Lambda(1-\bar s)},
 \qquad |\epsilon|=1.
\]
Suppose $\rho=1/2+i\gamma$ is a zero of order $m\ge1$.
Then, unconditionally,
\begin{equation}\label{eq:analytic-reflected-positive}
 \Re\frac{\Lambda'}{\Lambda}(\rho+1)
   =m+\sum_{\rho'\ne\rho}
       \frac{3/2-\Re\rho'}
       {(3/2-\Re\rho')^2+(\Im\rho'-\gamma)^2}>m.
\end{equation}
Multiplicity is retained and the displayed real series is
absolutely convergent.
\end{lemma}

\begin{proof}
Put $F(z)=\Lambda(\rho+z)$ and write its first two coefficients
as $b_mz^m+b_{m+1}z^{m+1}+O(z^{m+2})$. Since
$1-\bar\rho=\rho$, reflection becomes
$F(z)=\epsilon\overline{F(-\bar z)}$. Therefore
\[
 b_m=\epsilon(-1)^m\bar b_m,
 \qquad b_{m+1}=\epsilon(-1)^{m+1}\bar b_{m+1},
 \qquad
 \Re\frac{b_{m+1}}{b_m}=0.
\]
The genus-one Hadamard product has the form
\[
 F(z)=z^m e^{az+b}\prod_{\alpha\ne0}
           E_1(z/\alpha),\qquad
 E_1(w)=(1-w)e^w,
 \quad \alpha=\rho'-\rho.
\]
All multiplicities are included and
$\sum |\alpha|^{-2}<\infty$. Since each $E_1$ has zero
linear coefficient at the origin, the logarithmic derivative
of $F(z)/z^m$ at zero is $a=b_{m+1}/b_m$. Thus
$\Re a=0$; this is the required control of the exponential
factor, not an assumption about its sign.

At $z=1$, which is not a zero, logarithmic differentiation gives
\[
 \frac{F'(1)}{F(1)}=m+a+
     \sum_{\alpha\ne0}\left(\frac1{1-\alpha}
                                      +\frac1\alpha\right).
\]
The combined complex terms converge absolutely. Moreover
$|\Re\alpha|\le1/2$, so
$\sum|\Re(1/\alpha)|<\infty$. Reflection pairs $\alpha$ with
$-\bar\alpha$, with the same multiplicity, and
\[
 \Re\left(\frac1\alpha+\frac1{-\bar\alpha}\right)=0.
\]
For a purely imaginary zero this cancellation holds termwise.
Taking real parts therefore leaves just
$m+\sum\Re(1/(1-\alpha))$. Each summand is
\[
 \frac{1-\Re\alpha}
 {(1-\Re\alpha)^2+(\Im\alpha)^2}>0,
\]
because $1-\Re\alpha\ge1/2$. Its tail is bounded by a fixed
multiple of $\sum|\alpha|^{-2}$, proving absolute convergence.
There are noncentral zeros, so the inequality is strict.
Substituting $\alpha=\rho'-\rho$ proves
\eqref{eq:analytic-reflected-positive}.
\end{proof}

\begin{theorem}[NN-A035: strict left bending at multiple Dirichlet zeros]
\label{thm:analytic-multiple-dirichlet-bending}
Let $L(s)=L(s,\chi)$ for a nonprincipal primitive Dirichlet
character, real or complex, of parity $a\in\{0,1\}$ and
conductor $q$. Suppose $\rho=1/2+i\gamma$ is a zero of
multiplicity $m\ge2$. Write
\[
 L(\rho+z)=A_mz^m+A_{m+1}z^{m+1}+O(z^{m+2}),
 \qquad
 D_\rho=\frac1\rho+\frac{L'}L(\rho+1)
                              -\frac{A_{m+1}}{A_m}.
\]
Then
\begin{equation}\label{eq:analytic-multiple-D-positive}
 \Re D_\rho>m+h_a(\gamma)>m,
\end{equation}
where the exact positive functions are
\[
 h_1(\gamma)=\frac{\pi}{2\cosh(\pi\gamma)},\qquad
 h_0(\gamma)=\frac1{2(\gamma^2+1/4)}
                      -\frac{\pi}{2\cosh(\pi\gamma)}>0.
\]
For the shifted deformation of NN-A004, the $m$ normalized
branches after $t=u^m$ satisfy
\begin{equation}\label{eq:analytic-multiple-full-jets}
 s_j(t)=\rho+c_jt^{1/m}
                  +\frac{D_\rho}{m}c_j^2t^{2/m}
                  +O(t^{3/m}),\qquad
 c_j^m=\frac{\rho L(\rho+1)}{A_m}\ne0.
\end{equation}
Every branch has a nonzero real normal coefficient by its
second Puiseux order. Specifically, if $\Re c_j=0$ and
$c_j=i\beta_j$, then
\[
 \Re s_j(t)-\frac12
       =-\frac{\beta_j^2}{m}\Re D_\rho\,t^{2/m}
                         +O(t^{3/m})<0
\]
for all sufficiently small positive $t$.
If vertical first directions occur, the exact left/right
counts are $(m/2+1,m/2-1)$ for even $m$, and
$((m+1)/2,(m-1)/2)$ for odd $m$.
\end{theorem}

\begin{proof}
Use the exact completion
\[
 Q_\chi(s)=\left(\frac q\pi\right)^{(s+a)/2}
                         \Gamma((s+a)/2),
 \qquad \Lambda(s)=Q_\chi(s)L(s,\chi).
\]
The primitive functional equation and conjugation give
$\Lambda(s)=\epsilon_\chi
\overline{\Lambda(1-\bar s)}$ even when $\chi$ is complex;
the root number has modulus one. These are DLMF~25.15.5--25.15.6.
The theta integral as in NN-A032, now with the conjugate theta
series on the reflected half, gives the entire order-one
completion. The classical critical-strip and infinite-zero
inputs are DLMF~25.15(ii); cancellation at the endpoint uses
$L(1,\chi)\ne0$, DLMF~25.15.9, just as in NN-A032.
Thus Lemma~\ref{lem:analytic-reflected-logderivative} applies.

The coefficient ratio in the completed germ is
\[
 \frac{b_{m+1}}{b_m}
          =\frac{A_{m+1}}{A_m}+\frac{Q_\chi'}{Q_\chi}(\rho).
\]
Its real part is zero by reflection. Consequently
\begin{equation}\label{eq:analytic-multiple-D-completed}
 \Re D_\rho=
 \Re\frac{\Lambda'}\Lambda(\rho+1)
 +\Re\left(\frac1\rho+
       \frac{Q_\chi'}{Q_\chi}(\rho)
       -\frac{Q_\chi'}{Q_\chi}(\rho+1)\right).
\end{equation}
The conductor term cancels exactly. It remains to evaluate
the real gamma difference, retaining its height dependence.

Put $z=(\rho+a)/2$. Since $\Re\rho=1/2$, one has
$\overline{(\rho+1)/2}=1-\rho/2$. The digamma reflection
identity $\psi(1-z)-\psi(z)=\pi\cot(\pi z)$, DLMF~5.5.4,
therefore gives
\[
 \Re\left(\frac12\psi(\rho/2)
                   -\frac12\psi((\rho+1)/2)\right)
             =-\frac{\pi}{2\cosh(\pi\gamma)}.
\]
Indeed $\Re\cot(\pi/4+i\pi\gamma/2)=1/\cosh(\pi\gamma)$,
as follows by writing sine and cosine in exponentials.
This yields $h_0$ after adding
$\Re(1/\rho)=1/[2(\gamma^2+1/4)]$.
For odd parity, the recurrence (DLMF~5.5.2)
$\psi(1+\rho/2)=\psi(\rho/2)+2/\rho$ reverses the preceding
difference after the same $1/\rho$ term is added, giving
$h_1=\pi/(2\cosh(\pi\gamma))$.
The positivity of $h_0$ is exact: from $2<\pi<4$,
\[
 \cosh(\pi\gamma)\ge1+\frac{\pi^2\gamma^2}{2}
              >\pi(\gamma^2+1/4).
\]
This proves both positive functions. Combining them with
\eqref{eq:analytic-reflected-positive} proves
\eqref{eq:analytic-multiple-D-positive}.

For completeness, the local deformation coefficient is
derived rather than imported from a simple-zero formula.
The shifted expansion gives
\[
 L_t(\rho+z)=A_mz^m+A_{m+1}z^{m+1}
       -t(B_0+B_1z)+O(z^{m+2})+O(tz^2)+O(t^2),
\]
where $B_0=\rho L(\rho+1)$ and
$B_1=L(\rho+1)+\rho L'(\rho+1)$. Since $\Re(\rho+1)=3/2$,
the Euler product makes $B_0\ne0$. Put $t=u^m$ and
$z=c_ju+d_ju^2+O(u^3)$. The coefficient of $u^m$ gives
$A_mc_j^m=B_0$, with distinct nonzero roots. The coefficient
of $u^{m+1}$ gives
\[
 mA_mc_j^{m-1}d_j+A_{m+1}c_j^{m+1}-B_1c_j=0,
\]
because $2m\ge m+2$ for $m\ge2$. Thus
$d_j=c_j^2D_\rho/m$, proving
\eqref{eq:analytic-multiple-full-jets}. The analytic implicit
function theorem after division by $u^m$ justifies these
convergent sections and their common error estimate.

For a vertical $c_j=i\beta_j$, its second real coefficient
is strictly negative by \eqref{eq:analytic-multiple-D-positive}.
Every other branch is decided by its nonzero first real
coefficient. If $m$ is even, vertical directions come as the
opposite pair; the remaining roots are antipodal and split
equally between left and right. If $m$ is odd, there is exactly
one vertical root, and the other roots pair with opposite
real parts by their angular positions $\pi/2+2\pi j/m$ (or
the corresponding lower vertical starting angle). Adding
the left-bending vertical branches gives the stated counts.
\end{proof}

\begin{corollary}[Dedekind zeros at sufficient height, and every multiple zeta zero under A]
\label{cor:analytic-multiple-dedekind-bending}
Let $K$ have signature $(r_1,r_2)$ and degree $n$. At a
critical-line zero $\rho=1/2+i\gamma$ of $\zeta_K$ with
order $m\ge2$, the same $D_\rho$ and Puiseux expansion obey
\begin{equation}\label{eq:analytic-dedekind-D-lower}
 \Re D_\rho>m-
 \left(\frac{3}{2(\gamma^2+9/4)}
       +\frac{r_1\pi}{2\cosh(\pi\gamma)}
       +\frac{r_2}{2(\gamma^2+1/4)}\right).
\end{equation}
In particular, the sufficient condition
\[
 \gamma^2\ge\frac1m
       \left(\frac{r_2+3}{2}+\frac{r_1}{\pi}\right)
\]
already implies $\Re D_\rho>0$ and strict left bending of
every vertical first direction. If
\[
 \gamma^2\ge\frac{r_2+3}{2}+\frac{r_1}{\pi},
\]
then the stronger bound $\Re D_\rho>m-1>0$ holds. The simpler sufficient bound
$\gamma^2\ge n+2$ also works.

For $K=\mathbb Q$, the existing NN-A014 interval certificate
implies that every critical-line zero has $\gamma^2>8$.
Hence every multiple critical-line zeta zero satisfies the
stronger bound
\[
                    \Re D_\rho>m-\frac{11}{48}>0.
\]
Under A this treats all multiple nontrivial zeta zeros, without
assuming that any such zero exists.
\end{corollary}

\begin{proof}
Apply Lemma~\ref{lem:analytic-reflected-logderivative} to the
entire pole-removed Dedekind completion of NN-A033. The
local coefficient relation and
\eqref{eq:analytic-multiple-D-completed} still hold. The
pole-removal factors contribute
\[
 \Re\left(\frac1\rho+\frac1{\rho-1}
                                  -\frac1{\rho+1}\right)
                  =-\frac{3}{2(\gamma^2+9/4)}.
\]
Each real-place gamma factor contributes
$-\pi/(2\cosh(\pi\gamma))$, as just calculated. Each
complex-place factor contributes
$\Re(\psi(\rho)-\psi(\rho+1))=-\Re(1/\rho)$.
This proves \eqref{eq:analytic-dedekind-D-lower} with the
exact stated completion and no location hypothesis beyond
the specified zero being on the line.

For $\gamma\ne0$, the bracket is strictly less than
\[
 \frac1{\gamma^2}
        \left(\frac{r_2+3}{2}+\frac{r_1}{\pi}\right),
\]
using $\cosh(\pi\gamma)>\pi^2\gamma^2/2$. Comparing this
upper bound first with $m$ and then with $1$ proves both
thresholds. Since $\pi>3$ and $n=r_1+2r_2$,
the numerator is less than $n/3+3/2<n+2$, giving the
simpler sufficient threshold.

For the zeta assertion, reuse the positive theta integral
and exact normalization $H_0(x)=\xi(1/2+ix/2)/8$ of
NN-A006. The full saved Arb ball of NN-A014 is
\[
 A_0=\frac{\xi''(0)}{4\xi(1/2)}
       \in (0.0117395,0.0117396)\subset(0,1/16).
\]
No new computation is required for this wider rational
enclosure. If $c_1$ is the normalized second theta moment,
the positive coefficient sum in NN-A014 gives
$2c_1\le A_0<1/16$. The elementary inequality
$\cos v\ge1-v^2/2$ therefore gives, unconditionally,
\[
 H_0(x)\ge H_0(0)(1-c_1x^2)>0\qquad(x^2\le32).
\]
This is a statement about the integral and does not use a
real-zero product or RH. Thus a critical-line zero has
$x=2\gamma$ with $\gamma^2>8$. In the case $r_1=1,r_2=0$,
the bracket in \eqref{eq:analytic-dedekind-D-lower} is
strictly less than
$((3/2)+(1/3))/8=11/48$. The claimed lower bound and all
second-jet conclusions follow.
\end{proof}

\paragraph{Why the simple-zero gate remains open.}
For $m=1$, write $v=\rho L(\rho+1)/L'(\rho)$. Direct
coefficient matching instead gives
\begin{equation}\label{eq:analytic-simple-extra-term}
 s(t)=\rho+vt+
 \left(v^2D_\rho-
       \frac{v(\rho+1)L(\rho+2)}{2L(\rho+1)}\right)t^2
       +O(t^3).
\end{equation}
The extra term is the contribution of
$t^2\rho(\rho+1)L(\rho+2)/2$ in the shifted series.
If $v$ is purely imaginary, positivity of $\Re D_\rho$
controls only the first term in the parentheses. The real
part of the second term has not been controlled here.
Thus this theorem resolves the second normal jet at multiple
zeros, and neither proves $\Re v\ne0$ at simple zeros nor
assigns a universal bending direction to a simple tangent.

\paragraph{AG translation.}
At a length-$m$ critical-line analytic zero, completed
reflection puts the first reduced logarithmic coefficient
in $i\mathbb R$. The genus-one exponential and reflected
zero pairs make this local Real constraint compatible with
the strictly positive real logarithmic derivative one unit
to the right. After $t=u^m$, the finite analytic Cartier
divisor normalizes into $m$ sections. For $m\ge2$ the
specified shift has no $t^2$ contribution at the second
Puiseux order, so every vertical tangent acquires the
proved nonzero left normal coefficient. The local algebra
is $\mathbb C[z]/(z^m)$; its normalization and ramified
cover have the standard differentials/discriminant
interpretation of SGA~1, Expos\'e~I, section~3,
Proposition~1 and section~4, Proposition~10. This is a
second-jet theorem for the given analytic deformation,
not a purity statement or an arithmetic realization of
the shifted denominators.
